\section{Fault Tolerance and Recovery}
\label{sec:fault-tolerance}

Fault tolerance in this system relies on five mechanisms selected based on the type of fault.

\subsection{The failure model}
\label{sec:failure-model}

We assume crash-stop and network partition. Byzantine behaviour is out of scope. \Cref{tab:failures} lists what can fail
and what the system does about it.

\begin{table}[t]
\centering\small
\renewcommand{\arraystretch}{1.2}
\caption{What can fail, how it is noticed, and what it costs.}
\label{tab:failures}
\begin{tabularx}{\textwidth}{@{}>{\raggedright\arraybackslash}p{2.6cm}
                               >{\raggedright\arraybackslash}p{2.9cm}
                               X
                               >{\raggedright\arraybackslash}p{2.9cm}@{}}
\toprule
What fails & How it is noticed & What happens & What is lost \\
\midrule
Coordinator, crash & \texttt{nodedown}, at once &
  election; the new leader rebuilds from the stations &
  nothing \\
Coordinator, partition & heartbeat, 3~s &
  the minority suspends itself &
  new reservations, minority side \\
Station node, crash & \texttt{monitor\_node}, at once &
  the coordinator drops that station's claims &
  sessions running there \\
Station uplink, partition & \texttt{monitor\_node}, on the tick &
  the site keeps charging, starts nothing new &
  the operator's view \\
Charge point & socket closes, then 30~s grace &
  session closed with the last measured energy &
  nothing measured \\
Back office & nothing, it is a hidden node &
  the cluster is unaffected; it re-reads on restart &
  nothing durable \\
MySQL & connection errors &
  not tolerated, see \cref{sec:no-tolerance} &
  writes, until it returns \\
\bottomrule
\end{tabularx}
\end{table}

\subsection{Two failure detectors, because there are two failures}
\label{sec:detectors}

\texttt{net\_kernel:monitor\_nodes/1} comes for free and fires the instant a TCP
connection breaks, which is what a killed container looks like. A partition breaks
nothing: the peer is unreachable but its socket stays open, so no
\texttt{nodedown} is delivered until the distribution's own tick expires, and
\texttt{net\_ticktime} defaults to sixty seconds. A minority leader would go on
granting claims for a minute without knowing it had lost its quorum.

Liveness is therefore decided by a heartbeat among the coordinators, implemented
in \texttt{vs\_coord\_membership}: each announces itself once a second, and three
missed announcements mean gone. Three seconds instead of sixty. It runs among the
three coordinators only, because that is where a stale view produces two leaders;
a station that goes silent is still noticed on the tick, at the cost
\cref{tab:failures} states. \texttt{monitor\_nodes} stays subscribed for the crash
case, where it is faster still.

Measured on 1 September, isolating the leader with
\texttt{docker network disconnect}:

\begin{itemize}
  \item \texttt{QUORUM LOST (1 of 3) \ldots{} abdicating} at \textbf{2.12~s};
  \item a new leader serving with both claims at \textbf{2.29~s};
  \item the distribution's own \texttt{nodedown} at \textbf{64.6~s}.
\end{itemize}

That last number is the argument. Without the heartbeat the system would have
been incorrect for a full minute, and no test would have revealed it, because
\texttt{docker kill}, the obvious way to simulate a failure, closes the socket
and never exercises this path at all. Two kinds of failure need two experiments,
which is why the demonstration shows both \texttt{kill} and \texttt{disconnect}.

\subsection{A station dies}
\label{sec:station-failures}

\texttt{vs\_coord\_srv} monitors every station it has
heard from and, on \texttt{nodedown}, erases that station's claims. Keeping them
would shut those drivers out of the whole network until the lease expired.

The cars carry on charging, because the node is not on the path that delivers
power. The charge point is also the side that counts the energy, so when the
station comes back it reports its running total, the station adopts the session
from it, and the charge is billed in full. What is lost is a car that unplugs
while the node is still down: nothing writes its row, and that energy is never
invoiced.

\subsection{Rebuilding the claim table after a failover}
\label{sec:rebuild}

The new leader must discover the claims by querying the system.
It sends a \texttt{who\_do\_you\_hold} message to each station; the station examines its connectors and responds by indicating those in the \texttt{held} (occupied) or \texttt{charging} state, based on its local memory.
The leader transitions from the \texttt{rebuilding} state to the \texttt{serving} (operational) state once the responses are received.
Requests received in the interim are rejected with a retry indication, which the stations interpret as a temporary refusal.

\paragraph{Convergence via two paths.} 
The system would converge autonomously without the need for explicit queries from the new coordinator, as stations renew their claims every 10 seconds;
however, the query allows for a significant reduction in recovery time.

\subsection{A charge point dies}
\label{sec:cp-failure}

A lost socket is not a broken charger, so the connector waits ninety seconds
before doing anything: sixty of socket idle timeout, thirty of its own grace. If
they run out mid-session, the session is closed with the last energy the charge
point reported, the only figure anybody measured, and the connector goes
\texttt{out\_of\_service} until a charge point boots and reports itself available.

\subsection{Supervision inside a node}
\label{sec:supervision}


Three supervisory strategies have been implemented:

\begin{lstlisting}[caption={The two supervision trees. The station nests a second
                   supervisor inside the first; the coordinator does not.},
                   label={lst:suptrees}]
vs_station_sup                      one_for_one
  +-- vs_connector_sup              simple_one_for_one
  |     +-- vs_connector            connector 1   started at boot by
  |     +-- vs_connector            connector 2   vs_station_mgr, one
  |     +-- vs_connector            connector N   per physical connector
  +-- vs_station_mgr
  +-- vs_claim_client
  +-- vs_station_db

vs_coord_sup                        rest_for_one
  +-- vs_coord_srv                  claim table, cluster map
  +-- vs_coord_bo                   JInterface bridge to the back office
  +-- vs_coord_election             bully election
  +-- vs_coord_membership           heartbeat and quorum
      declaration order is dependency order
\end{lstlisting}

A supervisor restarts its children, and the strategy decides which ones go down
together. The two trees choose differently.

The coordinator is \texttt{rest\_for\_one}: its children are declared in
dependency order, so a process that dies takes with it everyone declared after
it. The bridge publishes what the claim table holds, and must not outlive a table
that was restarted underneath it.

The station is \texttt{one\_for\_one}: its children are independent and restart
alone. Connectors find the manager by registered name and hold no reference to
it, so a manager that comes back re-adopts them and no session is lost. The
connectors are all the same kind of process, one per outlet, hence
\texttt{simple\_one\_for\_one}; one that crashes comes back in \texttt{free}.

\subsection{Delivery semantics, chosen per message}
\label{sec:guarantees}

The bridge from the coordinator to the back office carries three kinds of event,
and they deliberately do not get the same guarantee, because the cost of a
duplicate is different for each.

\begin{itemize}
  \item \textbf{The session row: at-least-once over an idempotent write.} The
        message itself is never resent. What makes it at-least-once is that the
        station has already written the row at the end of the session, and the
        back office sweeps \texttt{cost\_cents IS NULL} every sixty seconds, so
        the event only makes an invoice arrive sooner and losing it costs one
        interval. The price is written with \texttt{UPDATE sessions SET
        cost\_cents = ? WHERE id = ? AND cost\_cents IS NULL}: a repeat updates
        zero rows, which is not an error.
  \item \textbf{The no-show strike: maybe.} One send, no retransmission and no
        duplicate filter. A counter cannot recognise a duplicate, so a retry
        could suspend an innocent driver, and a lost strike is a smaller wrong
        than an invented one.
  \item \textbf{The notification to the driver: maybe}, for the same reason. A duplicate is a
        row somebody reads twice, a loss is a row they never see, and we chose
        the quieter failure.
\end{itemize}

\subsection{What tolerates nothing, and is declared}
\label{sec:no-tolerance}

Three things here have no recovery path.

\textbf{MySQL is a single point of failure} for durable data. Replicating it was
out of reach for this project. This choice leads some limitations:

\textbf{Sessions on a dead station are lost}, as \cref{sec:station-failures}
describes. The process that was metering them is gone.

\textbf{The rebuild can still miss a station that is unreachable at exactly the
wrong moment.} Its claims are absent from the new leader's table, so one of its
vehicles could obtain a second claim before the conflict is noticed. The window
is narrow and it is real, and it closes when the station returns and re-presents
the older claim, which wins.
